\documentclass[12pt]{article}
\usepackage[T1]{fontenc}
\usepackage[polish]{babel}
\usepackage[utf8]{inputenc}
\usepackage[backend=biber,style=numeric,sortcites=true]{biblatex}
\usepackage{hyperref}
\usepackage{csquotes}
\usepackage{amsmath}
\usepackage{lmodern}
\usepackage{microtype}
\addbibresource{../Latex/bibliography_poczatek.bib}
\setlength{\emergencystretch}{3em} %~Allow LaTeX to stretch lines to avoid overfull boxes
\babelprovide[transforms =~oneletter.nobreak]{polish}
\begin{document}

%\section*{Pracować szybko, nie zatrzymywać się przy jednem problemie na dłużej niż to konieczne, potem będzie czas na poprawki}
%\begin{itemize}
%    \item [x] Zrobić reposytorium github i~projekt Overleaf
%    \item [x] Zrobić przegląd artykułów na temat
%    \item [~] Przygotować konspekt, cele i~aspekty
%    \item [x] Przerobić dane, żeby zorientować się co do ich jakości i~tego co zawierają
%    \item [x] Skontaktować się z~Panią Dyrektor i~byłym promotorem
%    \item [~] Wypełnić formularz WZI od nowa
%\end{itemize}

%\section*{Podsumowanie}

\section{Źródła danych}\label{sec:zroda-danych}

\subsection{Dane populacyjne}\label{subsec:dane-populacyjne}
\begin{itemize}
            \item Miejsce zamieszkania (ludność De Jure): Dane przestrzenne o~gęstości/rozmieszczeniu populacji dla
            wszystkich miast europejskich są szeroko dostępne. Najprzystępniejszą formą są dane z~badania Census,
            które zapisane są dla całego obszaru UE na poziomie granuralności siatki 1 km x~1 km. Problemem okazują
            się dane z~Australii. Tamtejszy odpowiednik Urzędu Statystycznego dzieli obszary w~niestandardowy sposób
            utrudniający analizę, dodatkowo sama możliwość eksportu tych danych jest mocno ograniczona. Problem ten
            ma przełożenie na wszystkie inne przestrzenne dane populacyjne dotyczące Australii.
            \item Miejsce zatrudnienia/populacja dzienna/populacja w~miejscu pracy (ludność De Facto): O~ile ten typ
            danych nie znajdował się początkowo w~obszarze moich poszukiwań, to może on jednak wykazywać wyższą
            rzeczywistą użyteczność w~analizie dostępności transportu miejskiego przy analizie uwzględniającej
            liczbę zatrudnionych. Dane takie dostępne są często wraz z~poprzednim typem danych, jednak dla Francji i~
            Niemiec nie zostały uwzględnione w~głównym spisie powszechnym, ale dostępne są do pobrania w~podobnej
            formie z~innych źródeł. Status dostępności danych dotyczących Australii taki sam jak wyżej.
            \item Dokładana lokalizacja miejsca pracy: Dane te różnią się od poprzednich tym, że określają dokąd
            ludzie udają się do pracy, a~nie ile osób na danym obszarze ma pracę. Dane tego typu pozwoliłyby mi
            uwzględnić potencjalne dzienne przepływy ludności, co znacząco podniosłoby trafność analiz, jednakże ich
            dostępność jest bardzo mocno ograniczona:
            \begin{itemize}
                \item Dublin -~zdecydowanie najlepsza dostępność danych, możliwe do pobrania są dane o~lokalizacji
                zatrudnienia w~siatce 1 km x~1km.
                \item Paryż -~pomimo istnienia danych takich, jak w~przypadku Dublina, nie są one publicznie
                dostępne (bądź nie udało mi się ich znaleźć, pomimo intesywnych poszukiwań), jednakże odnalazłem
                dane określające strefy o~wysokim zatrudnieniu, które pomimo trudności z~eksportem ich są możliwe do
                zescrappowania.
                \item Warszawa -~nie są dostępne publicznie dane w~stylu Dublina, czy nawet Paryża, jedynie udało mi
                się znaleźć dane określające liczbę ofert pracy w~podziale na dzielnice, co stanowi punkt wyjścia. O~
                ile się orientuję, to istnieje jeszcze możliwość zamówienia indywidualnej analizy poprzez GUS, ale nie
                zgłębiałem w~pełni tego tematu.
                \item Dla pozostałych miast, dla których szukałem danych tj. Lublin, Adelaide, Berlin, Brisbane,
                Canberra, Helsinki, Kuopio, Luxembourg, Melbourne, Rzym, Turku, Sydney, Winnipeg; nie udało mi się
                znaleźć tego typu danych.
            \end{itemize}
        \end{itemize}
    \subsection{Dane dotyczące komunikacji miejskiej}\label{subsec:dane-dotyczace-komunikacji-miejskiej}
\begin{itemize}
            \item Polska:

            Udało mi się znaleźć dane dotyczące wszystkich interesujących mnie przewoźników (ZTM Warszawa,
            Warszawska Kolej Dojazdowa, Koleje Mazowieckie, PolRegio), udostępnianie bezpośrednio przez nich i~
            skatalogowane na stronie \href{https://dane.gov.pl/pl/dataset/1739,krajowy-punkt-dostepowy-kpd-multimodalne-usugi-informacji-o-podrozach/resource/327976/table?page=1&per_page=20&q=warszaw&sort=}{dane.gov.pl}.
            \item Zbiór 25-cities:

            Informacje o~ofercie komunikacyjnej pozostałych miast pochodzić będą ze zbioru danych 25-cities
            \cite{25cities}, który zawiera dane GTFS dla 25 miast na świecie, w~tym Paryża i~Dublina. Nie wykluczam
            jednak możliwości pozyskania nowszych danych GTFS bezpośrednio od przewoźników lub z~oficjalnych
            repozytoriów danych otwartych.

        \end{itemize}

\section{Cele Pracy Magisterskiej}\label{sec:cele-pracy-magisterskiej}
Głównym celem pracy jest opracowanie i~zastosowanie holistycznych ram oceny dostępności transportu zbiorowego (TZ),
które wykraczają poza tradycyjne miary bazujące wyłącznie na podaży (ofercie przewozowej), w~celu przeprowadzenia
komparatywnej analizy trzech zróżnicowanych aglomeracji: Paryża, Dublina i~Warszawy.

Cel ten zostanie osiągnięty poprzez realizację następujących szczegółowych zadań:
\begin{enumerate}

    \item Rozszerzenie metodologiczne analizy dostępności (Paryż): Rozbudowanie wyników projektu zaliczeniowego o~
    zaawansowane miary dostępności, w~tym modele grawitacyjne oraz analizę topologiczną L-space (
    Space-of-Infrastructure) i~P-space (Space-of-Service). Użycie tych dualnych reprezentacji jest niezbędne do
    diagnozowania problemów sieciowych, rozróżniając, czy niski poziom dostępności wynika z~braku fizycznej
    infrastruktury (L-space), czy z~nieefektywnych operacji (P-space)~\cite{LUO2019102505}.

    \item Integracja Perspektywy Popytu i~Podaży (Modele Grawitacyjne): Kwantyfikacja potencjału dostępu do
    kluczowych celów (np. miejsc pracy, POI) z~wykorzystaniem modelu grawitacyjnego, najlepiej dwustronnie
    ograniczonego (doubly constrained). Model ten, poprzez zastosowanie współczynników normalizacyjnych ($A_i$ i~$B_i$),
    jest bardziej realistyczny niż model pojedynczo ograniczony, ponieważ uwzględnia konkurencję o~ograniczone
    zasoby (pojemność celów -~$E_j$) i~jednocześnie ogranicza całkowity wypływ z~obszaru początkowego (popyt -~$P_i$)
    ~\cite{gravity, gravity2}.

    \item Wzbogacenie Behawioralne Impedancji (Whole Travel Chain -~WTC): Uwzględnienie pełnego kosztu podróży z~
    perspektywy użytkownika, który jest kluczowy dla realistycznej miary impedancji $c_{ij}$. Oznacza to włącznie
    tzw. pierwszej i~ostatniej mili, szacunkowego czasu oczekiwania oraz czasu/penalizacji za przesiadki. Pominięcie
    tych składników prowadzi do przeszacowania użyteczności systemu TZ~\cite{LUO2019102505, capetown}.

    \item Weryfikacja Wymiaru Sprawiedliwości Społecznej (Mapowanie Społeczne): Przeprowadzenie analizy
    sprawiedliwości (equity), która, w~przeciwieństwie do równości (equality), wymaga dostosowanego wsparcia i~
    traktowania w~oparciu o~potrzeby grup. Należy to osiągnąć poprzez segmentację demograficzną modelu
    grawitacyjnego i~kalibrację współczynnika zaniku $\beta$ (wrażliwość na czas podróży) w~oparciu o~różnice w~
    zamożności (dochodach). Sprawiedliwe rezultaty w~transporcie publicznym (PT) są kluczowe dla redukcji
    nierówności, zwłaszcza w~dostępie do zatrudnienia, opieki zdrowotnej i~edukacji~
    \cite{KHALIFE2023160, capetown, urban_institute, review}.

    \item Analiza Komparatywna i~Zarządzanie Niespójnością Danych: Porównanie wyników dostępności uzyskanych za
    pomocą tej samej holistycznej metodologii dla trzech miast, z~uwzględnieniem, że analizy porównawcze są w~
    literaturze rzadkie, ponieważ istniejące metody są często zależne od GIS. Konieczne jest udokumentowanie
    horyzontów czasowych i~oszacowanie ryzyka zniekształcenia wyników z~powodu niespójności danych z~różnych okresów~\cite{LUO2019102505}

\end{enumerate}

\section{Aspekty Analityczne}\label{sec:aspekty-analityczne}
\begin{enumerate}

    \item Analiza Topologiczna (L-space i~P-space)
    \newline
    Zastosowanie topologii L-space i~P-space jest kluczowe w~analizie sieci transportu publicznego (PTN), ponieważ 
    różne reprezentacje topologiczne ujmują różne aspekty sieci~\cite{LUO2019102505, c-space}.
    \begin{itemize}

        \item Modelowanie Sieci Transportu Zbiorowego (PTN):

        \begin{itemize}
            \item L-space (Space-of-Infrastructure): Jest to bezpośrednia reprezentacja sieci (uproszczony graf), w~
            której wierzchołki reprezentują przystanki, a~krawędź łączy sąsiednie przystanki (jeśli na danej trasie
            istnieje co najmniej jedna usługa je obsługująca). L-space skupia się na fizycznym aspekcie
            infrastruktury~\cite{LUO2019102505, c-space, l-space2}.

            \item P-space (Space-of-Service): Jest to reprezentacja oparta wyłącznie na trasach (usługach).
            Wierzchołki również reprezentują przystanki, ale są ze sobą połączone krawędzią, jeśli są obsługiwane
            przez tę samą trasę (linia łączy wszystkie przystanki na danej trasie). P-space jest szeroko stosowany
            do badania możliwości przesiadek (transfer possibilities), ponieważ odzwierciedla bezpośrednią łączność,
            nie uwzględniając fizycznego układu infrastruktury (tj. sekwencji fizycznie pokonywanych odcinków)~\cite{LUO2019102505, c-space}.
        \end{itemize}

        \item Diagnostyka Problemów: Porównanie centralności lub miar dostępności obliczonych w~obu przestrzeniach
        pozwala zdiagnozować, czy niski poziom dostępności wynika z~braku infrastruktury (niski L-space) czy z~
        nieefektywnych operacji i~serwisu (niski P-space przy wysokim L-space, np. rzadkie kursy lub złe przesiadki)~\cite{LUO2019102505}.

        \item Źródła Danych: Wymaga to użycia danych GTFS dla każdego z~miast, które zawierają informacje o~trasach,
        przystankach i~sekwencjach~\cite{LUO2019102505, 25cities}.

    \end{itemize}

    \item Modele Grawitacyjne i~Impedancja
    \newline
    Modele grawitacyjne (Gravity-type approach)~\cite{gravity2} będą stanowiły rdzeń ilościowej oceny dostępności
    potencjalnej:

    \begin{itemize}

        \item Identyfikacja POI (Points of Interest): Wskaźnik atrakcyjności ($T_{attr}$ /~$O_j$) musi uwzględniać
        dzienne cele przepływu ludności, takie jak miejsca pracy (zatrudnienie), jako kluczowy element analizy
        dojazdów~\cite{gravity2}. Możliwe będzie wykorzystanie istniejących danych demograficznych (siatka 1x1 km)
        oraz specyficznych danych o~lokalizacji zatrudnienia, szczególnie dobrze opisanych dla Dublina.

        \item Aktualizacja metryki kosztu: Zamiast prostego syntetycznego wskaźnika, użytego w~projekcie zaliczeniowym,
        zastosowany zostanie uogólniony koszt podróży (Generalized Travel Cost -~GTC), który będzie stanowił
        podstawę funkcji impedancji $f(c_{ij})$. Koszt $c_{ij}$ powinien być oczekiwanym czasem podróży
        transportem zbiorowym~\cite{gravity2, LUO2019102505}.

        \item Współczynnik impedancji oparty o~zamożność ($\beta$): W~celu uwzględnienia sprawiedliwości społecznej
        (Society Mapping), współczynnik zaniku ($\beta$), określający wrażliwość na czas podróży, musi być
        kalibrowany w~oparciu o~różnice w~zamożności (dochody). Wysoki dochód często koreluje z~wyższą wrażliwością
        na czas (wyższe $\beta$), co wymaga segmentacji demograficznej modelu~\cite{gravity2}.

        \item (Potancjalnie) Dwustronnie ograniczony model grawitacyjny: Jest to najbardziej realistyczna forma
        modelu grawitacyjnego (zgodnie z~modelem Wilsona, 1969), która jednocześnie ogranicza przepływ z~obszaru
        początkowego (popyt) i~pojemność miejsca docelowego (podaż -~$D_j$/$e_j$)~\cite{gravity, gravity2}. W~
        przeciwieństwie do modeli kumulatywnych, modele grawitacyjne pozwalają na uwzględnienie efektów konkurencji
        o~ograniczone zasoby, co jest kluczowe dla realistycznej oceny~\cite{gravity2}.

    \end{itemize}

    \item Cały Łańcuch Podróży (Whole Travel Chain -~WTC)
    \newline
    Analiza WTC pozwoli na realistyczne uwzględnienie multimodalnego charakteru podróży. Koszt podróży $c_{ij}$
    powinien być postrzegany jako suma wszystkich segmentów.

    \begin{itemize}

        \item Elementy czasu podróży: Uogólniony koszt podróży (GTC) jest sumą całkowitego czasu jazdy pojazdem (
        in-vehicle time), całkowitego czasu oczekiwania na każdym etapie podróży (początkowe i~transferowe) oraz
        czasu/penalizacji za przesiadki~\cite{capetown, LUO2019102505}.

        \item Pierwsza i~Ostatnia Mila (First and Last Mile): Czas dostępu pieszo (do przystanku) i~dojścia od
        przystanku końcowego (first/last mile) są kluczowymi składnikami impedancji. W~przypadku braku szczegółowych
        danych (np. GTFS) czas oczekiwania ($T_{\text{wait}}$) może być oszacowany na podstawie częstotliwości
        kursowania (headway)~\cite{LUO2019102505, metropolitan}.

        \item Koszty Przesiadek: Czas przesiadki i~oczekiwania są postrzegane jako bardziej uciążliwe niż czas
        spędzony w~pojeździe. Analizy w~Nowej Zelandii wykazały, że użytkownicy z~niepełnosprawnościami potrzebują
        minimalnej oszczędności czasu, aby wybrać trasę z~przesiadką, a~poprawa jakości węzłów przesiadkowych nie
        miała pozytywnego wpływu na pożądany maksymalny czas oczekiwania i~chodzenia (sugerując, że duże węzły mogą
        być trudniejsze w~manewrowaniu)~\cite{disabilities1}.

        \item (Potencjalnie) Bariery Fizyczne i~Społeczne: Choć trudne do kwantyfikacji, należy uwzględnić, że
        bariery fizyczne (np. schody na stacjach metra) mogą znacznie obniżyć dostępność sieci dla osób z~
        niepełnosprawnościami, czyniąc dostępną sieć znacząco mniejszą. Ponadto, czynniki socjologiczne, takie jak
        postawa kierowców autobusów i~ich nieświadomość potrzeb pasażerów, są uznawane za znaczącą barierę (o wiele
        ważniejszą dla użytkowników niż dla projektantów)~\cite{elavating, disabilities1, disabilities3}.


%        \item Pierwsza i~Ostatnia Mila (First and Last Mile): Całkowity koszt podróży ($c_{ij}$) musi być obliczony
%        jako suma wszystkich segmentów: czas dostępu pieszo/rowerem (do przystanku), czas oczekiwania/transferu (w
%        tym penalizacja za przesiadki) oraz czas w~pojeździe~\cite{capetown}.
%        \item Źródła Danych WTC: Czas przejścia pieszego między przystankami jest już częściowo dostępny w~zbiorze
%        GTFS na podstawie OpenStreetMap~\cite{25cities}. Czas oczekiwania ($T_{wait}$) zostanie oszacowany na podstawie
%        częstotliwości kursowania (headway) z~danych GTFS~\cite{luo, metropolitan}. Pominięcie tych
%        kosztów (jak to często bywa w~uproszczonych modelach) prowadzi do przeszacowania użyteczności systemu TZ.
%        \item Włączenie Baterii (Bariery Fizyczne i~Społeczne): Choć trudne, bariery socjologiczne (np. w~przypadku
%        osób z~niepełnosprawnościami) mogą być włączone do modelu WTC jako negatywne stałe specyficzne dla
%        alternatywy ($ASC_m$), obniżające użyteczność całego łańcucha podróży dla wrażliwych grup.
    \end{itemize}

    \item Zarządzanie Niespójnością Czasową Danych
    \newline
    Fakt, że dane transportowe i~dane demograficzne/zatrudnienia pochodzą z~różnych okresów czasowych, może to
    prowadzić do zniekształcenia wyników~\cite{gravity} i~musi zostać jasno zaadresowany poprzez:

    \begin{itemize}

        \item Udokumentowanie horyzontów czasowych: Należy szczegółowo opisać, z~jakiego okresu pochodzi każda
        kategoria danych dla każdego miasta. Na przykład, w~zestawie \textit{25-cities} dane GTFS dla Dublina,
        Paryża i~innych miast pochodzą z~lat 2014-2016~\cite{25cities}.

        \item Strategię metodologiczną: Należy przyjąć spójne założenia, np. używać starszych danych populacyjnych
        jako statycznego rozkładu popytu dla analizy wrażliwości sieci transportowej z~danego okresu, a~niespójność
        ująć jako ograniczenie pracy

%        \item Szacowanie ryzyka zniekształcenia: Należy ocenić, w~jakim stopniu zmiany w~sieci transportowej (GTFS) lub
%        dystrybucji populacji/zatrudnienia mogły wpłynąć na wyniki w~okresie między pozyskaniem danych a~datą analizy.

    \end{itemize}
\end{enumerate}

%\subsection*{Konspekt Pracy Magisterskiej (Konspekt)}
%\begin{enumerate}
%    \item Wprowadzenie
%    \begin{enumerate}
%        \item Wstęp: Definicja dostępności obszarowej i~jej znaczenie w~planowaniu zrównoważonego transportu.
%        \item Kontekst Badawczy: Krótki przegląd istniejących miar dostępności (np. kumulatywne możliwości, modele potencjalne/grawitacyjne) \cite{czternascie}.
%        \item Ograniczenia Projektu Wstępnego: Krytyka wskaźnika użytego w~projekcie dla Paryża (skoncentrowanie na podaży/ofercie) i~uzasadnienie potrzeby przejścia na model grawitacyjny i~perspektywę behawioralną.
%        \item Cele i~Pytania Badawcze: Jasne sformułowanie celów pracy (patrz sekcja I) oraz hipotez badawczych.
%        \item Struktura Pracy.
%    \end{enumerate}
%\end{enumerate}

% Tytuł pracy –~powinien być krótki, ale jednocześnie precyzyjnie odzwierciedlać tematykę pracy.
% Wstęp –~wprowadzenie do tematyki pracy, cel pracy oraz pytania badawcze, na które praca ma odpowiedzieć.
% Stan wiedzy –~przegląd literatury na temat badanej problematyki, z~uwzględnieniem najważniejszych publikacji w~danym zakresie.
% Metodologia –~opis wykorzystanych metod badawczych, źródeł i~narzędzi badawczych.
% Plan badań –~opis planowanych badań, w~tym opis próby badawczej, sposobu zbierania danych oraz analizy wyników.
% Wyniki badań –~prezentacja wyników badań oraz ich analiza.
% Dyskusja –~interpretacja wyników badań oraz wnioski wyciągnięte z~przeprowadzonych badań.
% Podsumowanie –~podsumowanie najważniejszych wyników badań oraz wskazanie kierunków dalszych badań w~danym obszarze.

%\subsection*{Konspekt Pracy Magisterskiej}
%\begin{enumerate}
%    \item Tytuł pracy
%
%    ,,Holistyczna analiza przestrzennej dostępności transportu zbiorowego w~aglomeracjach miejskich''
%    \item Wprowadzenie
%    \item Źródła danych
%        \begin{itemize}
%        \item Populacyjne
%            \begin{itemize}
%                \item Miejsce zamieszkania (ludność De Jure): Dane przestrzenne o~gęstości/rozmieszczeniu populacji dla wszystkich miast europejskich są szeroko dostępne. Najprzystępniejszą formą są dane z~badania Census, które zapisane są dla całego obszaru UE na poziomie granuralności siatki 1 km x~1 km. Problemem okazują się dane z~Australii. Tamtejszy odpowiednik Urzędu Statystycznego dzieli obszary w~niestandardowy sposób utrudniający analizę, dodatkowo sama możliwość eksportu tych danych jest mocno ograniczona. Problem ten ma przełożenie na wszystkie inne przestrzenne dane populacyjne dotyczące Australii.
%                \item Miejsce zatrudnienia/populacja dzienna/populacja w~miejscu pracy (ludność De Facto): O~ile ten typ danych nie znajdował się początkowo w~obszarze moich poszukiwań, to może on jednak wykazywać wyższą rzeczywistą użyteczność w~analizie dostępności transportu miejskiego przy analizie uwzględniającej liczbę zatrudnionych. Dane takie dostępne są często wraz z~poprzednim typem danych, jednak dla Francji i~Niemiec nie zostały uwzględnione w~głównym spisie powszechnym, ale dostępne są do pobrania w~podobnej formie z~innych źródeł. Status dostępności danych dotyczących Australii taki sam jak wyżej.
%                \item Dokładana lokalizacja miejsca pracy: Dane te różnią się od poprzednich tym, że określają dokąd ludzie udają się do pracy, a~nie ile osób na danym obszarze ma pracę. Dane tego typu pozwoliłyby mi uwzględnić potencjalne dzienne przepływy ludności, co znacząco podniosłoby trafność analiz, jednakże ich dostępność jest bardzo mocno ograniczona:
%                \begin{itemize}
%                    \item Dublin -~zdecydowanie najlepsza sytuacja, gdzie dostępne są dane o~lokalizacji zatrudnienia w~siatce 1 km x~1km.
%                    \item Paryż -~pomimo istnienia danych takick, jak w~przypadku Dublina, nie są one publicznie dostępne (bądź nie udało mi się ich znaleźć, pomimo intesywnych poszukiwań), jednakże odnalazłem dane określające strefy o~wysokim zatrudnieniu, które pomimo trudności z~eksportem ich są możliwe do zescrappowania.
%                    \item Warszawa -~nie są publicznie dostępne dane w~stylu Dublina, czy nawet Paryża, jedynie udało mi się znaleźć dane określające liczbę ofert pracy w~podziale na dzielnice, co stanowi punkt wyjścia. O~ile się orientuję, to istnieje jeszcze możliwość zamówienia indywidualnej analizy poprzez GUS, ale zgłębiałem w~pełni tego tematu.
%                    \item Dla pozostałych miast, dla których szukałem danych tj. Lublin, Adelaide, Berlin, Brisbane, Canberra, Helsinki, Kuopio, Luxembourg, Melbourne, Rzym, Turku, Sydney, Winnipeg; nie udało mi się znaleźć tego typu danych.
%                \end{itemize}
%            \end{itemize}
%        \item Dotyczące komunikacji miejskiej
%            \begin{itemize}
%                \item Polska:
%
%                Udało mi się znaleźć dane dotyczące wszystkich interesujących mnie przewoźników (ZTM Warszawa, Warszawska Kolej Dojazdowa, Koleje Mazowieckie, PolRegio), udostępnianie bezpośrednio przez nich i~skatalogowane na stronie \href{https://dane.gov.pl/pl/dataset/1739,krajowy-punkt-dostepowy-kpd-multimodalne-usugi-informacji-o-podrozach/resource/327976/table?page=1&per_page=20&q=warszaw&sort=}{dane.gov.pl}.
%                \item Zbiór 25-cities:
%
%                Informacje o~ofercie komunikacyjnej będę dalej czerapał głównie ze zbioru danych
%
%            \end{itemize}
%        \end{itemize}
%    \item Metodologia
%    \item Struktura pracy
%
%        \begin{minipage}{\linewidth}
%        % \textbf{\large Mini Table of Contents}\par\vspace{0.7em}
%
%        \noindent 1. Wstęp \dotfill \par
%        \noindent 2. Wstęp \dotfill \par
%        \noindent \hspace{1.5em} 1.1 Mini Subsection A \dotfill \par
%        \noindent \hspace{1.5em} 1.2 Mini Subsection B \dotfill \par
%        \noindent 2. Mini Section Two \dotfill \par
%        \end{minipage}
%
%    \item Podsumowanie
%\end{enumerate}

\section{Konspekt Pracy Magisterskiej}\label{sec:konspekt-pracy-magisterskiej}

\subsection*{Kontekst badawczy}

    W~literaturze dostępności transportu publicznego dominują miary oparte głównie na
    podaży (np. liczba kursów, zasięg sieci) lub proste miary kumulatywne, które nie zawsze odzwierciedlają
    rzeczywiste możliwości dotarcia do miejsc pracy i~usług. Jest potrzeba opracowania holistycznych ram oceny
    dostępności, łączących reprezentacje topologiczne sieci (L-space i~P-space), modele grawitacyjne uwzględniające
    konkurencję o~zasoby oraz behawioralnie ugruntowaną miarę impedancji (Whole Travel Chain), a~także analizę
    sprawiedliwości społecznej (equity). Równocześnie porównawcze analizy między miastami są rzadkie ze względu na
    niespójność źródeł danych (czasowych i~formatowych) i~różną dostępność danych GTFS czy lokalizacji miejsc pracy.
    Praca ma wypełnić tę lukę poprzez zastosowanie spójnej metodologii porównawczej dla Paryża, Dublina i~Warszawy.

\subsection*{Pytania badawcze}

    \begin{enumerate}

        \item W~jaki sposób holistyczne ramy łączące L-space/P-space, dwustronnie ograniczone modele grawitacyjne i~
        Whole Travel Chain zmieniają ocenę dostępności transportu publicznego w~porównaniu do klasycznych miar
        opartych wyłącznie na podaży, na przykładach Paryża, Dublina i~Warszawy?

        \item Jakie różnice w~diagnozie problemów sieciowych (infrastruktura vs. serwis) ujawniają porównania miar
        centralności/dostępności obliczonych w~L-space i~P-space dla każdego z~miast?

        \item W~jakim stopniu model grawitacyjny (w porównaniu z~modelami kumulatywnymi)
        poprawia odzwierciedlenie rzeczywistego potencjału dostępu do miejsc pracy i~innych POI?

        \item Jak uwzględnienie first/last mile, czasu oczekiwania i~kosztów przesiadek (GTC) wpływa na wartości
        funkcji impedancji i~w konsekwencji na szacunki dostępności w~badanych miastach?

        \item Czy i~jak różni się kalibracja współczynnika zaniku $\beta$ między grupami o~różnym statusie
        dochodowym -- tzn. czy osoby o~niższych dochodach wykazują inną wrażliwość na czas podróży i~jak to wpływa na
        sprawiedliwość dostępu?

        \item Które z~trzech miast (Paryż, Dublin, Warszawa) wykazują największe ograniczenia dostępności wynikające z:
        \begin{itemize}
            \item braku infrastruktury,
            \item słabego serwisu (częstotliwości/przesiadek),
            \item niekorzystnej struktury geograficznej miejsc pracy (rozproszenie/koncentracja)?
        \end{itemize}

        \item Jakie rekomendacje administracyjne wynikają z~analizy (np. poprawa częstotliwości, inwestycje w~węzły
        przesiadkowe, działania pierwszej/ostatniej mili) w~kontekście zwiększania dostępu do
        zatrudnienia i~usług?

        \item Jak niespójność czasowa i~różna dostępność danych (np. GTFS, lokalizacji miejsc pracy) wpływa na
        porównania między miastami i~jak zastosowane metody minimalizują to ryzyko zniekształcenia wyników?

    \end{enumerate}

\section{Przykładowa struktura Pracy}\label{sec:przykadowa-struktura-pracy}

\subsection*{Rozdział 1: Przegląd Literatury i~Podstawy Metodologiczne}
\begin{enumerate}
    \item Teoretyczne ramy dostępności: cztery komponenty dostępności —~transportowy, przestrzenny, czasowy,
    indywidualny~~\cite{capetown}.
    \item Topologiczna analiza sieci (L-space i~P-space): definicja, różnice i~znaczenie dualizmu L-space/P-space w~
    diagnozowaniu problemów infrastrukturalnych i~operacyjnych~\cite{LUO2019102505}.
    \item Modele grawitacyjne w~ocenie dostępności: podstawy modelu, rola atrakcyjności celów ($T_{\text{attr}}$ --
    POI, zatrudnienie) i~impedancji ($f(c_{ij})$)~\cite{gravity, gravity2}.
    \item Modelowanie całego łańcucha podróży (WTC): konieczność włączenia pierwszej i~ostatniej mili (access/egress).
    Włączenie kosztów oczekiwania i~transferów (penalizacja) do funkcji GTC, która jest podstawą impedancji~\cite{LUO2019102505}.
    \item Sprawiedliwość transportowa i~mapowanie społeczne: koncepcja równości (equality) i~sprawiedliwości (
    justice). Rola segmentacji demograficznej w~kalibracji współczynnika zaniku $\beta$ (wrażliwość na czas a~
    zamożność) oraz uwzględnienie, że bariery pozamaterialne (np. fizyczne przeszkody) mają kluczowe znaczenie dla
    osób z~niepełnosprawnościami~\cite{disabilities31, disabilities3, elavating}.
\end{enumerate}

\subsection*{Rozdział 2: Dane i~Rozwiązanie Problemu Niespójności Czasowej}
\begin{enumerate}
%    \item Studia przypadku: uzasadnienie wyboru Paryża (kontynuacja), Dublina (dostępność danych przepływów) i~Warszawy (miasto lokalne, polski kontekst).
    \item Dostępność danych transportowych (GTFS):
    \begin{itemize}
        \item Źródła GTFS dla Paryża i~Dublina (dane z~zestawu\textit{25-cities} z~2016 r.\ i~2018 r.).
        \item Wyzwania i~rola GTFS w~modelowaniu sieci (czasy podróży, częstotliwości, itp.).
    \end{itemize}
    \item Dostępność danych popytowych i~celów (POI/zatrudnienie):
    \begin{itemize}
        \item Ludność (de jure): dane z~badania Census.
        \item Wskaźniki zamożności: dane z~badania Census (jeśli to nie wystarczy to opisać konieczność pozyskania lub
        użycia zastępczych wskaźników społeczno-ekonomicznych dla segmentacji)~\cite{KHALIFE2023160, elavating}.
        \item Paryż/Dublin/Warszawa: omówienie dostępnych danych o~zatrudnieniu.
    \end{itemize}
    \item Adresowanie niespójności czasowej danych:
    \begin{itemize}
        \item GTFS z~różnych lat (np. Paryż 2016) oraz dane demograficzne/zatrudnienia z~innych okresów.
        \item Strategie normalizacji/harmonizacji: omówienie sposobów korekty różnic czasowych między danymi
        transportowymi a~demograficznymi. Ujęcie tego jako ograniczenia.
    \end{itemize}
\end{enumerate}

\subsection*{Rozdział 3: Metodologia Zaawansowanej Analizy Dostępności}
\begin{enumerate}
    \item Etap I: przygotowanie danych i~topologia sieci:
    \begin{itemize}
        \item Integracja i~czyszczenie danych GTFS (autobusy, metro/kolej, tramwaje) i~danych OSM dla pieszych.
        \item Modelowanie sieci dla każdego miasta w~dualizmie L-space (infrastruktura) i~P-space (usługa).
        \item Analiza podstawowych właściwości topologicznych sieci (centralność, spójność).
    \end{itemize}
    \item Etap II: kalibracja kosztu podróży (WTC):
    \begin{itemize}
        \item Obliczenie uogólnionego kosztu podróży (GTC) $c_{ij}$, obejmującego cały łańcuch podróży.
        \item Włączenie elementów pierwszej i~ostatniej mili: czas dotarcia pieszego (OSM), czas oczekiwania (
        częstotliwość —~$T_{\text{wait}}$) oraz czas transferu (penalizacja).
    \end{itemize}
    \item Etap III: implementacja modeli grawitacyjnych:
    \begin{itemize}
        \item Definicja stref analitycznych (np.\ siatka 1×1 km).
        \item Identyfikacja i~kwantyfikacja celów (POI) oraz potencjału zatrudnienia ($T_{\text{attr}}$).
        \item Zastosowanie modelu grawitacyjnego (dwustronnie ograniczonego) do obliczenia macierzy przepływów.
        \item Kalibracja współczynnika zaniku $\beta$ dla różnych grup społeczno-ekonomicznych.
        \item Obliczenie wskaźnika dostępności $A_i$ (np.\ odwrócony średni koszt dojazdu).
    \end{itemize}
    \item Etap IV: mapowanie społeczne i~porównanie:
    \begin{itemize}
        \item Wizualizacja i~analiza przestrzenna nierówności w~dostępie dla różnych grup społecznych.
        \item Porównanie miar dostępności między Paryżem, Dublinem i~Warszawą.
    \end{itemize}
\end{enumerate}

\subsection*{Rozdział 4: Wyniki i~Dyskusja}
\begin{enumerate}
    \item Wyniki topologiczne (L/P-space): diagnoza problemów infrastruktury i~operacyjnych w~każdym mieście.
    \item Wyniki grawitacyjne i~WTC: prezentacja map dostępności potencjalnej i~wpływu uwzględnienia WTC oraz
    segmentacji $\beta$ na wyniki, w~kontekście redukcji lub zaostrzenia spatial disparity w~dostępności.
    \item Analiza sprawiedliwości transportowej: kwantyfikacja różnic w~dostępie (np.\ do miejsc pracy) w~zależności
    od zamożności strefy/grupy docelowej.
    \item Dyskusja komparatywna i~ograniczenia: analiza wpływu różnic w~jakości/okresie danych na porównywalność miast.
\end{enumerate}

\subsection*{Wnioski i~Rekomendacje}
\begin{enumerate}
    \item Synteza osiągniętych celów.
    \item Rekomendacje planistyczne: ukierunkowane interwencje dla każdego miasta (np.\ wzmocnienie L-space, poprawa
    częstotliwości, wsparcie mikromobilności -- „ostatnia mila”).
    \item (?) Kierunki dalszych badań.
\end{enumerate}

\newpage
\begin{raggedright}
\nocite{*}
\printbibliography
\end{raggedright}
\end{document}
